% RC Ladeschaltung mit Wechselschalter erweitertes Labeling
\begin{tikzpicture}
    \ctikzset{bipoles/cuteswitch/thickness=0.3}
    % RC Netzwerk mit Schalter und Wechselspannungsquelle uq. 
    % vor t0: Kurzschluss R2, parallel dazu Reihe aus R3 und C.
    % nach t0: lineare Spannungsquelle (uq mit R1) an R2, parallel dazu Reihe aus R3 und C.

    %  *---R1---o..........
    %  |..>t0.->..............
    %  |..........S1---*--R3--*
    % uq..<t0.->./.....|......|
    %  |........o......R2.....C
    %  |........|......|......|
    %  *--------*------*------*

    %  *---R1---o
    %  | t>t0 ->
    %  |          S1---*--R3--*
    % uq t<t0 -> /     |      |
    %  |        o      R2     C
    %  |        |      |      |
    %  *--------*------*------*

    %Switch Circuit [Position der nodes von Schalter] 
    % [s1: (2, 1.69); unten s1.cout 1: (1.70,?); oben s1.out 2: (?, 2.11); rechts s1.in: (?, 1.685)]  
    \draw(2,1.69)
        node[cute spdt up arrow, rotate = 180](s1){}% Cute Switch, Position des Schalters mit up/down veränderbar
        (s1.center) node[left, xshift= -5]{$t_0$};%Label 
        
    \draw(s1.out 2) to[R, l=$R_1$, name=R1, v_>, !vi] (0,2.11);% nach Uq+
    \draw([yshift=-0.5mm] s1.cout 1) to[short,-*] (1.70,0);%nach Uq-
    %RC Circuit
    \draw(0,2.11) 
    to[V, name=Uq, v_>, !v](0,0)
    to[short, -*](2.75,0)
    to[short, *-](5,0) 
    to[C, l=$C$, name=C, v<, !v](5,1.685)
    to[R, l=$R_3$, name=R3, v<, i^<, -*, !vi](2.75,1.685)
    to[R, l_=$R_2$, name=R2, v^, !vi](2.75,0);

    \draw(2.75,1.685) to[short](s1.in); % R2+ bis Schalter rechts

    % Strom- und Spannungspfeile mit Label
    \varrmore{Uq}{$u_q$};
    \varrmore{C}{$u_C$};
    \varrmore{R2}{$u_2$};

\end{tikzpicture} 
